\chapter[Comparison with Existing Frameworks]{Comparison with Existing Frameworks}\label{ch:sec:comparison-with-existing-frameworks}

This chapter situates the theory of weighted incidence structures,
their forced Hilbert-space geometry, and the Universal Optimality
Defect within the broader landscape of cryptography, information
theory, and geometric data analysis.  The goal is not to argue that
other frameworks are wrong, but to clarify what is recovered,
what is generalised, and what is genuinely new.

\section{Shannon's Perfect Secrecy}\label{ch:sec:shannon-perfect-secrecy}

Shannon's definition of perfect secrecy requires that the ciphertext
be independent of the message: $\Pr[M=m\mid C=c]=\Pr[M=m]$.
In the language of weighted incidence structures, this says that the
posterior equals the prior on every cloak.  With a uniform prior this
is equivalent to the vanishing of the Universal Optimality Defect
$\mathcal D(P)$ for every posterior distribution $P$; more generally,
the defect vanishes precisely when the posterior is proportional to
the prior, the condition of universal optimality (Section~\ref{ch:sec:wis}).

\textbf{What we recover.}
Perfect secrecy is the zero of our geometric scale: $\mathcal D\equiv0$.
The AES S-box, our main running example, satisfies this condition
exactly with a uniform prior: the incidence is complete, $n_{mc}=1$ for
every pair, and by Bayes' rule the posterior equals the uniform prior
on every cloak.

\textbf{What we generalise.}
Shannon's theory provides a binary classification (perfect secrecy
or not).  Our framework replaces the binary with a continuous measure:
the UOD quantifies \emph{how far} a system is from perfect secrecy,
not just whether it achieves it.

\textbf{What is new.}
The geometric structure---the Riemannian metric, the Lipschitz bounds,
the Three-Level Theorem---exists only when the defect is non-zero.
Shannon's framework, being purely probabilistic, has no access to this
geometry.

\section{Universal Optimality}\label{ch:sec:comparison-universal-optimality}

The concept of universal optimality was introduced by
Biondi~et~al.~\cite{biondi2018apollonian} to unify several security
notions under a single algebraic condition.  A deterministic encoder
is universally optimal if, after observing the ciphertext, every
compatible message is equally likely---a condition that coincides with
perfect secrecy when the key is uniform.

\textbf{What we recover.}
The characterisation of universal optimality as the condition that the
induced multiplicities factor as $n_{mc}=\beta(c)/w(m)$ is restated
and proved in this book (Theorem~\ref{ch:thm:representation}).

\textbf{What we generalise.}
The original theory treated universal optimality as an algebraic
property of multiplicities.  Our framework embeds it in a continuous
geometry, showing that every encoder---optimal or not---defines a point
in a Riemannian manifold whose distance from the optimal set is the
UOD.

\textbf{What is new.}
The geometric bounds (Lipschitz, quadratic, Bregman), the stability
theory, and the Three-Level variational characterisation are entirely
new and have no counterpart in the algebraic theory.

\section{Information Geometry and the Fisher Metric}\label{ch:sec:comparison-fisher}

Classical information geometry~\cite{amari2016information} studies the
space of probability distributions endowed with the Fisher--Rao metric.
The posterior manifold constructed in this book is a submanifold of the
simplex, but its metric is the pullback of the Euclidean metric through
the logarithmic map, not the Fisher metric.

\textbf{What we recover.}
Both geometries share the same underlying manifold---the interior of
the simplex---and both produce quadratic approximations for
information-theoretic quantities.  In the limit of small
perturbations, the pullback metric and the Fisher metric coincide up
to a factor of~2, because $\log(1+\delta)\approx\delta$ and the
Fisher metric is the pullback of the Euclidean metric through the
square-root map.

\textbf{What is different.}
The Fisher metric is invariant under sufficient statistics and arises
from the Cram\'er--Rao lower bound.  Our pullback metric arises from
the algebraic structure of the WIS factorisation, not from estimation
theory.  The two metrics differ globally: the UOD is a squared norm
in a Hilbert space, while the Fisher--Rao distance is geodesic
distance on the sphere (in the square-root representation).

\textbf{What is new.}
Our geometry is not chosen by the analyst; it is forced by the
representation condition $n_{mc}=\beta(c)/w(m)$.  This gives every
geometric statement---every bound, every stability result---a
cryptographic interpretation that information geometry alone cannot
provide.

\section{Optimal Transport and Wasserstein Distance}\label{ch:sec:comparison-optimal-transport}

Optimal transport theory~\cite{villani2009optimal} defines distances
between probability distributions as the minimum cost of moving mass
from one to another.  The Wasserstein distance $W_p(P,Q)$ has become
a standard tool in machine learning and statistics.

\textbf{What is different.}
The Wasserstein distance measures the \emph{effort} of transforming
one distribution into another, while the UOD measures the
\emph{distance from a structural condition}.  They are fundamentally
different objects: the UOD is a function of a single distribution
$\mathcal D(P)$ (or of a pair $(P,Q)$ through the defect), while
the Wasserstein distance is always a function of two distributions.

\textbf{What we recover.}
In the special case where $P$ is the uniform distribution,
$\mathcal D(P)=\|P-U_n\|^2$ recovers the $L^2$ distance, which is
related to the Wasserstein-2 distance through the Poincar\'e
inequality on the simplex.

\textbf{What is new.}
The UOD satisfies Lipschitz bounds with respect to the manifold metric,
not the Wasserstein metric.  Our bounds are stronger because the
manifold metric is adapted to the cryptographic structure, not to the
geometry of optimal transport.

\section{Graph Laplacians and Spectral Graph Theory}\label{ch:sec:comparison-graph-laplacians}

The posterior Laplacian $L=A^{\mathsf T}QA=D_M-HD_C^{-1}H^{\mathsf T}$
is the Schur complement~\cite{zhang2005schur} of the normal operator associated with the
incidence graph.  This operator is a generalised Laplacian on a
bipartite graph with vertex weights.

\textbf{What we recover.}
For an unweighted, regular bipartite graph, $L$ reduces to the
combinatorial Laplacian of the graph, whose spectral properties are
well studied~\cite{chung1997spectral}.  The least eigenvalue of $L$
determines the optimal Lipschitz constant for quadratic functionals.

\textbf{What generalises.}
Classical spectral graph theory studies a fixed graph with uniform
weights.  Our Laplacian depends on the posterior distribution $P$
through the weights $\pi_i$ and $P_i$, making it a function on the
manifold rather than a fixed operator.  This distribution-dependent
Laplacian is what produces the non-trivial geometry of the posterior
manifold.

\textbf{What is new.}
The connection between the normal operator, the Schur complement,
and the posterior Laplacian reveals that the UOD is the reduced
energy of a least-squares problem on the incidence graph---a
perspective that has no analogue in spectral graph theory.

\section{Convex Analysis and Bregman Divergences}\label{ch:sec:comparison-bregman}

Bregman divergences~\cite{bregman1967relaxation} are defined by
$D_\phi(x,y)=\phi(x)-\phi(y)-\langle\nabla\phi(y),x-y\rangle$ for
a strictly convex function $\phi$.  They are widely used in
optimisation, information theory, and machine learning.

\textbf{What we recover.}
The Average Hessian Operator (Chapter~\ref{ch:sec:the-average-hessian-operator-and-universal-bregman-theo})
recovers the exact integral representation
$D_\phi(P,Q)=\frac{1}{2}\Delta^{\mathsf T}\overline H_\phi\Delta$,
which is the Bregman identity expressed through the manifold geometry.

\textbf{What is new.}
The UOD bounds the spectral norm of the Average Hessian uniformly
across all smooth functionals.  This gives a universal estimate for
the curvature of arbitrary Bregman divergences on the posterior
manifold---a result that does not follow from classical convex
analysis because it depends on the specific geometry of the WIS
factorisation.

\section{Summary: What Is New}\label{ch:sec:comparison-summary}

The following table summarises the relationship between existing
frameworks and the theory developed in this book.

\begin{center}
\begin{tabular}{lll}
\hline
\textbf{Framework} & \textbf{Recovered} & \textbf{New} \\
\hline
Shannon perfect secrecy & $\mathcal D\equiv0$ & Continuous measure, geometry \\
Universal optimality & $n_{mc}=\beta/w$ & Riemannian structure, bounds \\
Information geometry & Simplex, quadratic approx & Forced metric, cryptographic meaning \\
Optimal transport & $L^2$ distance special case & Structural distance, Lipschitz bounds \\
Spectral graph theory & Combinatorial Laplacian & Distribution-dependent operator \\
Bregman divergences & Integral representation & Universal spectral bounds \\
\hline
\end{tabular}
\end{center}
